\solchapter{condmat}

\sol{condmat}{1} $\sin\theta = n\lambda/2d$: $0.273$ for $n=1$ ($\theta = 15.8^\circ$) and $0.546$ for $n = 2$ ($33.1^\circ$). Visible light has $\lambda\approx500\un{nm}\gg2d = 0.56\un{nm}$, so $\sin\theta>1$: no Bragg reflection is possible.

\sol{condmat}{2} $\tau = m\sigma/ne^2 = (9.11\times10^{-31})(6\times10^7)/[(8.5\times10^{28})(1.6\times10^{-19})^2] = 2.5\times10^{-14}\un{s}$. The mean free path is $v_F\tau = 40\un{nm}$, about $160$ atomic spacings: electrons pass many ions without scattering, as Bloch's theorem says they should.

\sol{condmat}{3} Visible photons span $1.7$--$3.1\un{eV}$. \emph{Silicon} ($1.1\un{eV}$) absorbs all of them and looks grey and opaque; an LED would emit infrared near $1.1\un{\mu m}$ (and a poor one, since silicon's gap is indirect). \emph{Gallium nitride} ($3.4\un{eV}$) is transparent; it emits near-ultraviolet, and alloyed with indium it gives the blue LEDs behind white LED lighting (Nobel Prize, 2014). \emph{Cadmium sulfide} ($2.4\un{eV}$) absorbs violet and blue and looks yellow-orange; it would emit green. \emph{Diamond} ($5.5\un{eV}$) is transparent, and would emit deep ultraviolet.

\sol{condmat}{4} The band runs from $E_0-2t$ to $E_0+2t$: width $4t$. $v = \hbar^{-1}\dd E/\dd k = (2ta/\hbar)\sin ka$. Near $k = 0$, $E\approx E_0 - 2t + ta^2k^2$, which matches $\hbar^2k^2/2m^*$ for $m^* = \hbar^2/2ta^2$. At the zone boundary $k = \pi/a$, $v = 0$: Bragg reflection turns the electron into a standing wave. Near the top of the band, the curvature is negative, so the effective mass is negative: an electron there accelerates opposite to the applied force. That is why it is simpler to describe the few empty states near the top as positively charged holes with positive mass.

\sol{condmat}{5} Modes with $\abs{\vb k}<k_T = \kB T/\hbar v$ fill a sphere in $k$-space containing $V\cdot\tfrac43\pi k_T^3/(2\pi)^3$ wavevectors, times three polarizations: a number $\propto VT^3$. Each holds energy $\sim\kB T$, so $U\propto VT^4$ and $C = \dd U/\dd T\propto T^3$. Photons are counted identically (two polarizations, speed $c$), so their energy density grows as $T^4$, and so does the flux they radiate: the Stefan--Boltzmann law. A cold crystal is a box of thermal sound.

\sol{condmat}{6} $\exp[(E_g/2\kB)(1/300 - 1/350)] = \exp[(6380\un{K})(4.8\times10^{-4}\un{K^{-1}})]\approx e^{3.0}\approx20$ (about $26$ including the $T^{3/2}$ prefactor). Heat increases the number of intrinsic carriers exponentially, swamping the designed doping and multiplying leakage currents, which heat the device further. Doping fixes the carrier density at a value set by the number of dopant atoms, nearly independent of temperature, which is what makes devices predictable.

\sol{condmat}{7} $\mu_BB/\kB T = 5.8\times10^{-5}/0.0259 = 2.2\times10^{-3}$, and the net alignment is $\tanh(\mu_BB/\kB T)\approx0.2\%$. Alignment reaches $50\%$ when $\tanh x = 0.5$, i.e.\ $x = 0.55$, at $T = 1.2\un{K}$. Nuclear moments are about $1000$ times smaller, so the nuclear polarization in an MRI scanner is only a few parts per million.

\sol{condmat}{8} $h/2e = 2.07\times10^{-15}\un{Wb}$. A thousandth of this through $1\un{mm^2}$ is $2\times10^{-12}\un{T}$: about $10^7$ times smaller than the Earth's field, but still twenty times larger than the brain's fields. Magnetoencephalography uses SQUIDs with larger pickup coils, longer averaging, and magnetically shielded rooms to reach the femtotesla level.

\sol{condmat}{9} $\lambda_L = \sqrt{m^*/\mu_0n^*e^{*2}} = \sqrt{2m_e/(\mu_0\cdot10^{28}\cdot4e^2)}\approx40\un{nm}$. A magnetic field dies away within a few tens of nanometres of a superconductor's surface.

\sol{condmat}{10} $V_H = IB/(net) = (10)(1)/[(8.5\times10^{28})(1.6\times10^{-19})(10^{-4})]\approx7\un{\mu V}$. In a semiconductor the carrier density is a million to a hundred million times lower, so the Hall voltage is millivolts to volts: easy to measure, and the basis of the Hall sensors in phones and cars.

\sol{condmat}{11} Label each cell's two atoms $A$ and $B$, with hopping $t_1$ between $A$ and $B$ in the same cell and $t_2$ between $B$ and the next cell's $A$. Bloch states $(a,b)e^{ikjd}$, with cell length $d = 2a$, satisfy $Ea = -(t_1 + t_2e^{-ikd})b$ and $Eb = -(t_1 + t_2e^{ikd})a$. Multiplying, $E^2 = \abs{t_1+t_2e^{ikd}}^2 = t_1^2 + t_2^2 + 2t_1t_2\cos kd$, so $E = \pm\abs{t_1+t_2e^{ikd}}$. The smallest magnitude, $\abs{t_1-t_2}$, occurs at $kd = \pi$, leaving a gap of $2\abs{t_1-t_2}$ between the bands. With one electron per atom the lower band is exactly full: the chain is an insulator. In one dimension, the electronic energy gained by opening the gap always outweighs the elastic cost of the distortion, so a half-filled chain dimerizes spontaneously.

\sol{condmat}{12} Threading flux $\Phi$ through the cylinder induces an EMF $\mathcal{E} = -\dd\Phi/\dd t$ around it, which drives a Hall current $I = \sigma_{xy}\mathcal{E}$ along the cylinder, from one edge to the other. Adding exactly $\Phi_0 = h/e$ changes every electron's wavefunction by a phase that can be removed by a gauge transformation, so the Hamiltonian and its spectrum return to their original form. If the bulk has an energy gap and the flux is added slowly, the system returns to its ground state, except that some charge may have been moved from edge to edge --- and since electrons are indivisible, it must be an integer number $\nu$ of them. So $\nu e = \int I\,\dd t = \sigma_{xy}\int\mathcal{E}\,\dd t = \sigma_{xy}h/e$, and $\sigma_{xy} = \nu e^2/h$. The quantization comes from the integrality of charge and the gap; disorder and sample shape never enter.
